% GENERATED FILE. Edit canonical lessons, then run npm run generate:books.
% canonical-source-hash: fnv1a64:767ab5c76b477dad
% canonical-lessons: JA-C215-karada, JA-C215-yubi, JA-C215-ude, JA-C215-kubi, JA-C215-seki, JA-R215-review-210-215

\chapter{Arms, Fingers and Necks}
\label{ch:ja-arms-fingers-and-necks}
\hlchaptermodality{215}

\begin{chapteropening}
\textbf{By the end of this chapter:} \emph{I can ask and answer about parts of the body.}

Complete the last lesson of chapter 215: I can ask and answer about parts of the body.
\end{chapteropening}

\section[\texorpdfstring{\ja{からだ} (karada)}{からだ (karada)}]{\ja{からだ} (karada) — the body}

\label{lesson:JA-C215-karada}



\begin{quote}
Before the new one: say the Japanese for to burn, then the Japanese for to vanish.
\end{quote}

\subsection*{You'll want to know: \ja{からだ}}
\textbf{\ja{からだ}} — \emph{karada} — \textquotedblleft{}the body\textquotedblright{}.

\begin{quote}
\textbf{\ja{これは} \ja{からだに} \ja{いいですか。}} — \emph{kore wa karada ni ii desu ka} — \textquotedblleft{}Is this good for your body?\textquotedblright{}
\end{quote}

\begin{cousinweb}
\textbf{\ja{からだ}} (\emph{karada}) is a native Japanese word (wago), inherited rather than borrowed from Chinese; its kanji was attached later as a kun'yomi, and its older origin is not settled (JMdict, KANJIDIC2).
\end{cousinweb}

\begin{grammarlens}[title={What you've built}]
One more word for everyday life at A2.
\end{grammarlens}

\subsection*{Your turn}
\begin{itemize}
\raggedright
  \item \emph{Say it:} \emph{karada}
  \item \emph{Say it:} \emph{karada}, once more
  \item \emph{Say it:} \emph{kore wa karada ni ii desu ka} — Is this good for your body?
  \item \emph{Recall:} say the Japanese for to burn, then the Japanese for to vanish, then say \emph{karada} again
\end{itemize}

\begin{tcolorbox}[breakable,colback=teal!4,colframe=teal!35!black,title={Before you move on}]
What does \ja{からだ} mean? (\textquotedblleft{}The body\textquotedblright{}.) Say it once more.
\end{tcolorbox}
\section[\texorpdfstring{\ja{ゆび} (yubi)}{ゆび (yubi)}]{\ja{ゆび} (yubi) — a finger; a toe}

\label{lesson:JA-C215-yubi}



\begin{quote}
Before the new one: say the Japanese for to vanish, then the Japanese for the body.
\end{quote}

\subsection*{You'll want to know: \ja{ゆび}}
\textbf{\ja{ゆび}} — \emph{yubi} — \textquotedblleft{}a finger; a toe\textquotedblright{}.

\begin{quote}
\textbf{\ja{ゆびが} \ja{うごきません。}} — \emph{yubi ga ugokimasen} — \textquotedblleft{}My finger does not move.\textquotedblright{}
\end{quote}

\begin{cousinweb}
\textbf{\ja{ゆび}} (\emph{yubi}) is a native Japanese word (wago), inherited rather than borrowed from Chinese; its kanji was attached later as a kun'yomi, and its older origin is not settled (JMdict, KANJIDIC2).
\end{cousinweb}

\begin{grammarlens}[title={What you've built}]
One more word for everyday life at A2.
\end{grammarlens}

\subsection*{Your turn}
\begin{itemize}
\raggedright
  \item \emph{Say it:} \emph{yubi}
  \item \emph{Say it:} \emph{yubi}, once more
  \item \emph{Say it:} \emph{yubi ga ugokimasen} — My finger does not move.
  \item \emph{Recall:} say the Japanese for to vanish, then the Japanese for the body, then say \emph{yubi} again
\end{itemize}

\begin{tcolorbox}[breakable,colback=teal!4,colframe=teal!35!black,title={Before you move on}]
What does \ja{ゆび} mean? (\textquotedblleft{}A finger; a toe\textquotedblright{}.) Say it once more.
\end{tcolorbox}
\section[\texorpdfstring{\ja{うで} (ude)}{うで (ude)}]{\ja{うで} (ude) — an arm}

\label{lesson:JA-C215-ude}



\begin{quote}
Before the new one: say the Japanese for the body, then the Japanese for a finger.
\end{quote}

\subsection*{You'll want to know: \ja{うで}}
\textbf{\ja{うで}} — \emph{ude} — \textquotedblleft{}an arm\textquotedblright{}.

\begin{quote}
\textbf{\ja{ちちは} \ja{うでが} \ja{つよいです。}} — \emph{chichi wa ude ga tsuyoi desu} — \textquotedblleft{}My father has strong arms.\textquotedblright{}
\end{quote}

\begin{cousinweb}
\textbf{\ja{うで}} (\emph{ude}) is a native Japanese word (wago), inherited rather than borrowed from Chinese; its kanji was attached later as a kun'yomi, and its older origin is not settled (JMdict, KANJIDIC2).
\end{cousinweb}

\begin{grammarlens}[title={What you've built}]
One more word for everyday life at A2.
\end{grammarlens}

\subsection*{Your turn}
\begin{itemize}
\raggedright
  \item \emph{Say it:} \emph{ude}
  \item \emph{Say it:} \emph{ude}, once more
  \item \emph{Say it:} \emph{chichi wa ude ga tsuyoi desu} — My father has strong arms.
  \item \emph{Recall:} say the Japanese for the body, then the Japanese for a finger, then say \emph{ude} again
\end{itemize}

\begin{tcolorbox}[breakable,colback=teal!4,colframe=teal!35!black,title={Before you move on}]
What does \ja{うで} mean? (\textquotedblleft{}An arm\textquotedblright{}.) Say it once more.
\end{tcolorbox}
\section[\texorpdfstring{\ja{くび} (kubi)}{くび (kubi)}]{\ja{くび} (kubi) — a neck}

\label{lesson:JA-C215-kubi}



\begin{quote}
Before the new one: say the Japanese for a finger, then the Japanese for an arm.
\end{quote}

\subsection*{You'll want to know: \ja{くび}}
\textbf{\ja{くび}} — \emph{kubi} — \textquotedblleft{}a neck\textquotedblright{}.

\begin{quote}
\textbf{\ja{くびが} \ja{いたいですか。}} — \emph{kubi ga itai desu ka} — \textquotedblleft{}Does your neck hurt?\textquotedblright{}
\end{quote}

\begin{cousinweb}
\textbf{\ja{くび}} (\emph{kubi}) is a native Japanese word (wago), inherited rather than borrowed from Chinese; its kanji was attached later as a kun'yomi, and its older origin is not settled (JMdict, KANJIDIC2).
\end{cousinweb}

\begin{grammarlens}[title={What you've built}]
One more word for everyday life at A2.
\end{grammarlens}

\subsection*{Your turn}
\begin{itemize}
\raggedright
  \item \emph{Say it:} \emph{kubi}
  \item \emph{Say it:} \emph{kubi}, once more
  \item \emph{Say it:} \emph{kubi ga itai desu ka} — Does your neck hurt?
  \item \emph{Recall:} say the Japanese for a finger, then the Japanese for an arm, then say \emph{kubi} again
\end{itemize}

\begin{tcolorbox}[breakable,colback=teal!4,colframe=teal!35!black,title={Before you move on}]
What does \ja{くび} mean? (\textquotedblleft{}A neck\textquotedblright{}.) Say it once more.
\end{tcolorbox}
\section[\texorpdfstring{\ja{せき} (seki)}{せき (seki)}]{\ja{せき} (seki) — a cough}

\label{lesson:JA-C215-seki}



\begin{quote}
Before the new one: say the Japanese for an arm, then the Japanese for a neck.
\end{quote}

\subsection*{You'll want to know: \ja{せき}}
\textbf{\ja{せき}} — \emph{seki} — \textquotedblleft{}a cough\textquotedblright{}.

\begin{quote}
\textbf{\ja{せきは} \ja{でません。}} — \emph{seki wa demasen} — \textquotedblleft{}I do not have a cough.\textquotedblright{}
\end{quote}

\begin{cousinweb}
\textbf{\ja{せき}} (\emph{seki}) is a native Japanese word (wago), inherited rather than borrowed from Chinese; its kanji spelling came later, and its older origin is not settled (JMdict).
\end{cousinweb}

\begin{grammarlens}[title={What you've built}]
One more word for everyday life at A2.
\end{grammarlens}

\subsection*{Your turn}
\begin{itemize}
\raggedright
  \item \emph{Say it:} \emph{seki}
  \item \emph{Say it:} \emph{seki}, once more
  \item \emph{Say it:} \emph{seki wa demasen} — I do not have a cough.
  \item \emph{Recall:} say the Japanese for an arm, then the Japanese for a neck, then say \emph{seki} again
\end{itemize}

\begin{tcolorbox}[breakable,colback=teal!4,colframe=teal!35!black,title={Before you move on}]
What does \ja{せき} mean? (\textquotedblleft{}A cough\textquotedblright{}.) Say it once more.
\end{tcolorbox}
\section[(dialogue)]{Words from chapters 210–215 — a review}

\label{lesson:JA-R215-review-210-215}



\begin{quote}
Say the Japanese for a tree and to plant.
\end{quote}

\subsection*{Your turn}
Say these aloud:

\begin{itemize}
\raggedright
  \item a tree, to plant, a plant, to grow, to wither
  \item an animal, a monkey, a rabbit, a small bird, to increase
  \item air temperature, to rise, to get cold, to freeze, to melt
  \item to shine, the moon, to sink, to climb, to get warm
  \item a storm, to shake, a fire, to burn, to vanish
  \item the body, a finger, an arm, a neck, a cough
\end{itemize}

\begin{tcolorbox}[breakable,colback=teal!4,colframe=teal!35!black,title={Before you move on}]
A tree? (\textbf{\ja{き}}, \emph{ki}.) To plant? (\textbf{\ja{うえる}}, \emph{ueru}.) A plant? (\textbf{\ja{しょくぶつ}}, \emph{shokubutsu}.) To grow? (\textbf{\ja{そだつ}}, \emph{sodatsu}.) To wither? (\textbf{\ja{かれる}}, \emph{kareru}.) An animal? (\textbf{\ja{どうぶつ}}, \emph{dōbutsu}.) A monkey? (\textbf{\ja{さる}}, \emph{saru}.) A rabbit? (\textbf{\ja{うさぎ}}, \emph{usagi}.) A small bird? (\textbf{\ja{ことり}}, \emph{kotori}.) To increase? (\textbf{\ja{ふえる}}, \emph{fueru}.) Air temperature? (\textbf{\ja{きおん}}, \emph{kion}.) To rise? (\textbf{\ja{あがる}}, \emph{agaru}.) To get cold? (\textbf{\ja{ひえる}}, \emph{hieru}.) To freeze? (\textbf{\ja{こおる}}, \emph{kōru}.) To melt? (\textbf{\ja{とける}}, \emph{tokeru}.) To shine? (\textbf{\ja{ひかる}}, \emph{hikaru}.) The moon? (\textbf{\ja{つき}}, \emph{tsuki}.) To sink? (\textbf{\ja{しずむ}}, \emph{shizumu}.) To climb? (\textbf{\ja{のぼる}}, \emph{noboru}.) To get warm? (\textbf{\ja{あたたまる}}, \emph{atatamaru}.) A storm? (\textbf{\ja{あらし}}, \emph{arashi}.) To shake? (\textbf{\ja{ゆれる}}, \emph{yureru}.) A fire? (\textbf{\ja{ひ}}, \emph{hi}.) To burn? (\textbf{\ja{もえる}}, \emph{moeru}.) To vanish? (\textbf{\ja{きえる}}, \emph{kieru}.) The body? (\textbf{\ja{からだ}}, \emph{karada}.) A finger? (\textbf{\ja{ゆび}}, \emph{yubi}.) An arm? (\textbf{\ja{うで}}, \emph{ude}.) A neck? (\textbf{\ja{くび}}, \emph{kubi}.) A cough? (\textbf{\ja{せき}}, \emph{seki}.)
\end{tcolorbox}
